\section{Prime-field research: sparse sets, root finding, and algebraic bases}\label{sec:prime}

Prime fields are the most demanding generalisation target in this
portfolio. Coordinate arithmetic is efficient, but the subfield
conditions used by extension-field IC are absent. The programme tests
several substitutes rather than treating this absence as a proof that no
alternative can exist.\src{PRIME}

\subsection{What the smallness experiments actually learned}

The dedicated smallness repository asks whether an inexpensive
coordinate-defined set behaves additively better than a matched random
set. Its experiments include intervals in $x$ and $y$, slopes, affine and
M\"obius images, Hamming weight, Farey heights, multiplicative subgroups,
random windows, Fourier spectra, low-degree polynomial sweeps,
Coppersmith constructions, Witt carries and learned log-bit predictors.
The main closure panel uses three roughly 20-bit prime-order toy curves.
\src{SMALL}

An illustrative observation is especially clear. For a base near size
2,048, ordinary coordinate predicates have pooled two-sum closure ratios
near one: $x<2048$ gives about 1.04, a subgroup predicate about 0.95, and
matched random sets about 0.98--1.00. A log-interval positive control,
which uses unavailable discrete-log information to define the set, gives
256.4 against a predicted 256.5. The control shows that the instrument
can see strong additive structure; the measured cheap predicates do not
show that structure on the tested panel.

The null comparison must match sign closure as well as cardinality.
Negation-closed sets have correlated pair sums, so a naive Poisson error
bar understates their spread. Random-window maxima must also be compared
with maxima from the same-size null search. The lesson is a concrete
experimental discipline, not a universal theorem that every inexpensive
predicate is random.

The Witt-carry and small-root studies examine why tempting analogues of
extension coordinates become expensive. Their polynomial degrees and
success thresholds belong to the specific constructions tested. A
failed learned log-bit predictor is similarly evidence about the models,
features, labels and curve sizes used. It does not prove that no
computational representation of prime-field structure can help.

\subsection{A useful reduction of the two-summand oracle}

For a short-Weierstrass model $y^2=x^3+ax+b$ in characteristic other
than two or three, and affine points $R,F$ with
$x_R\ne x_F$, the familiar chord formulas yield
\[
 S_3(x_R,x_F,X)=(x_R-x_F)^2
       \bigl(X-x(R+F)\bigr)\bigl(X-x(R-F)\bigr).
\]
Thus, once $F$ is fixed, the roots needed for two-summand decomposition
are already group-law expressions. One can subtract $F$ and test base
membership, handling signs and exceptional points explicitly, rather than
compute a fresh Gr\"obner basis. The proof source formalises the affine
algebraic identity; its coordinate definitions and hypotheses are visible
in the Lean file.\src{RF-PROOF}

This is an implementation opportunity, not a new general summation-law
theorem or a prime-field ECDLP breakthrough. The source audit explicitly
identifies prior sum-and-membership algorithms. It also corrects the
manuscript's broad lower-bound language: a reusable pair table can make a
repeated query constant-time after quadratic setup, so an unqualified
linear-per-query lower bound would be false.\src{RF-AUDIT}

The implementation family includes batched inversions for direct prime
subtraction, binary pair-sum tables, three- and four-summand meet-in-the-middle
variants, windowed scans, dense root indexes, Kummer-base variants and a
CUDA root-finding kernel. The prime and binary models use their own group
laws. The short-Weierstrass identity should not be silently substituted
for the generalized binary model $y^2+xy=x^3+ax^2+b$.\src{RF-PDP}
\src{RF-GPU}

\subsection{Measured improvements, with the comparator repaired}

The source audit reports binary per-decomposition ratios of 383 and
1,535 against matrix F4 at field degrees nine and thirteen. The
precomputed pair-table build is excluded from those query ratios. At
degree 23 a complete binary solve is reported 122 times faster with the
pair table, while F4 is faster at degrees five, nine and thirteen. This
reversal is evidence that setup and scale matter.\src{RF-AUDIT}

On five NIST P-curves the six-point-base query ratios are
3,021--5,802 against the benchmark's Buchberger implementation. The printed
benchmark label said ``F4'', but the timed source calls
\texttt{groebner\_f4::buchberger}. The corrected comparator belongs in
every reuse of those ratios. These experiments perform tiny-base oracle
work on large field parameters; they do not solve an unknown NIST-curve
logarithm.

The complete prime-DLP comparison against msolve has two evidence tiers.
At 16 bits, both totals are directly measured. At 20 and 24 bits, the F4
totals are composed from sampled per-query CPU cost and the root-finder's
measured query count. The reported factors of about 100, 105 and 50 must
not all be described as directly timed complete solves. Nor are any of
them comparisons against the strongest generic rho method.

The ecbench audit supplies a different view. On one registered 20-bit
prime subgroup, a pair table reduces the counted IC $S$ from 1,524.511 to
28.753 relative to direct subtraction. Its paired counted ratio is
0.018861. Native work remains unpriced on those prime arms, so that
quotient does not bound a complete-cost ratio. On a degree-17 Koblitz
instance, the corresponding IC comparison is fully priced in that
study's calibration: $356.894\to20.360$, paired ratio 0.057047. Even
there, the rho comparison has unpriced work and cannot be promoted into
a complete-cost speed claim.\src{RF-AUDIT}

\subsection{Kummer towers and the remaining algebraic question}

Petit--Kosters--Messeng-style bases use multiplicative and tower
structure to give summation systems more algebraic organisation.
\cite{pkm2016} In the recorded five-curve ecbench root-finding panel,
896 runs verify and the IC/rho ratios remain about 29--141. Changing a
plain base to a tower base of the same size changes that table-based
oracle's cost by roughly one percent. The conclusion is narrow: the
tower's algebraic structure is not exploited by that particular
enumeration/table mechanism. It says nothing decisive about a genuinely
algebraic tower solver.\src{PKM}

The large-memory work now addresses that distinction. At the pinned main
snapshot, several larger tower cells remain bounded or resource-limited.
An open branch subsequently records the $m=3,N=21$ run on a 124-GB,
32-vCPU host. Both targets finish all degree-seven reductions without a
refutation, then stop at the first degree-eight step under the frozen
dense-cap rule. The demanded $B'$ has 79,330,205,871 entries, approximately
317 GB at four bytes per entry; the recorded peak resident memory is
67.1 GB. The two runs take 6,445 and 6,416 seconds.\src{PR2160}

This is real execution progress relative to an unrun memory-gated card.
It does not supply a completed degree or a productive degree-eight
refutation. The branch's results explicitly leave the interpretation
between the protocol's inconclusive and lower-bound cases to the
coordinator. It would be inaccurate either to keep calling the cell
``never run'' or to call it a successful degree-eight solve.

Another open branch integrates prime tower F4, signature F5, an F6-style
closure and enumeration into ecbench. It reports 140 verified runs on
four 9--15-bit prime-order curves. F5 reduces the emitted F4 work on the
same trees, while F6's closures save essentially no total work. The
author-reported priced comparisons still favour enumeration and rho.
The first attempt is explicitly voided for thread-dependent replay.
These are unmerged engineering observations at their pinned PR head,
not an update silently applied to main's frontier.\src{PR2223}

\subsection{Earlier prime-curve structural work}

The May--August record also includes $p$-adic lifts, anomalous-curve
techniques applied outside their usual setting, CM analysis, GLV/HNP
lattices, twists and coordinate correlations. Its value includes
identifying which proposed bridge was absent. An efficient scalar
decomposition, a low-height coordinate or a numerical lattice anomaly
does not automatically yield an efficient unknown-scalar recovery
algorithm. The later LLL correction makes this especially tangible:
apparent differences between named curves can be implementation
artefacts.\src{HNP}\src{LLL}\src{STRUCTURAL}
